\documentclass[11pt,a4paper]{article}
\usepackage[T1]{fontenc}
\usepackage[utf8]{inputenc}
\usepackage{lmodern}
\usepackage[margin=27mm,headheight=15pt]{geometry}
\usepackage{amsmath,amssymb,amsthm,mathtools,bm}
\usepackage{microtype}
\usepackage{xcolor}
\usepackage{booktabs,longtable,array}
\usepackage{enumitem}
\usepackage{fancyhdr}
\usepackage{xurl}
\usepackage{hyperref}
\usepackage{bookmark}
\definecolor{ink}{HTML}{17324D}
\definecolor{linkblue}{HTML}{245A81}
\hypersetup{colorlinks=true,linkcolor=linkblue,citecolor=linkblue,urlcolor=linkblue,
 pdftitle={Gamma and Zeta over the Surcomplex Numbers},
 pdfauthor={AI-assisted research study prepared for Vladimir Reshetnikov}}
\urlstyle{same}
\setlength{\emergencystretch}{3em}
\setlength{\parskip}{4pt}
\setlength{\parindent}{15pt}
\setlist{nosep,leftmargin=1.8em}
\pagestyle{fancy}
\fancyhf{}
\fancyhead[L]{\small\textsc{Gamma and zeta over the surcomplex numbers}}
\fancyhead[R]{\small September 2026}
\fancyfoot[C]{\thepage}
\renewcommand{\headrulewidth}{0.3pt}
\numberwithin{equation}{section}
\newtheorem{theorem}{Theorem}[section]
\newtheorem{proposition}[theorem]{Proposition}
\newtheorem{corollary}[theorem]{Corollary}
\newtheorem{lemma}[theorem]{Lemma}
\theoremstyle{definition}
\newtheorem{definition}[theorem]{Definition}
\newtheorem{example}[theorem]{Example}
\theoremstyle{remark}
\newtheorem{remark}[theorem]{Remark}
\newtheorem{warning}[theorem]{Warning}
\newcolumntype{P}[1]{>{\raggedright\arraybackslash}p{#1}}
\newcommand{\No}{\mathbf{No}}
\newcommand{\Oz}{\mathbf{Oz}}
\newcommand{\SC}{\mathbb K}
\newcommand{\R}{\mathbb R}
\newcommand{\C}{\mathbb C}
\newcommand{\N}{\mathbb N}
\newcommand{\Z}{\mathbb Z}
\newcommand{\OO}{\mathcal O}
\newcommand{\mm}{\mathfrak m}
\newcommand{\st}{\operatorname{st}}
\newcommand{\pp}{\operatorname{pp}}
\newcommand{\fin}{\operatorname{fin}}
\newcommand{\cis}{\operatorname{cis}}
\newcommand{\Arg}{\operatorname{Arg}}
\newcommand{\Log}{\operatorname{Log}}
\newcommand{\Mer}{\operatorname{Mer}}
\newcommand{\ord}{\operatorname{ord}}
\newcommand{\RH}{\mathsf{RH}}
\newcommand{\Sh}{\mathop{\sum}\nolimits^{\mathrm H}}
\newcommand{\Ph}{\mathop{\prod}\nolimits^{\mathrm H}}
\newcommand{\DD}{\mathcal D}
\newcommand{\LSt}{L_{\mathrm{St}}}
\newcommand{\zD}{\zeta_{\mathrm D}}
\newcommand{\zp}{\zeta^{\#}}
\newcommand{\Gp}{\Gamma^{\#}}
\newcommand{\xip}{\xi^{\#}}
\newcommand{\val}{v}
\newcommand{\repo}{\texttt{VladimirReshetnikov/Surreal}}
\newcommand{\pin}{\texttt{905dd19954db43ac81f76f72036520d0eadc5710}}
\newcommand{\smallO}{\mathrm{o}}
\newcommand{\bigO}{\mathrm{O}}
\newcommand{\symh}{\mathcal S}
\newcommand{\ds}{\displaystyle}

\begin{document}
\begin{titlepage}
\vspace*{16mm}
{\color{ink}\Huge\bfseries Gamma and Zeta over the\\[4pt] Surcomplex Numbers\par}
\vspace{10mm}
{\Large Canonical finite lifts, infinite Dirichlet series,\\[3pt]
phase obstructions, and the Riemann hypothesis\par}
\vspace{13mm}
{\large An AI-assisted research study prepared for\\[3pt]
Vladimir Reshetnikov\par}
\vspace{5mm}
September 2026
\vspace{14mm}
\begin{abstract}
There is more than one mathematically defensible meaning of a Gamma or zeta
function with surcomplex arguments. We construct the canonical Taylor--Laurent
extensions on the finite surcomplex plane and prove exact preservation of their
divisors. Their Riemann hypothesis is therefore equivalent to the ordinary one.
A stronger stability statement is proved: all finite zeros of every symmetric
infinitesimal perturbation of the completed zeta function lie on the critical
line if and only if the ordinary Riemann hypothesis holds and all its
nontrivial zeros are simple.

At positive infinite real part, ordinary Dirichlet coefficients of arbitrary
growth admit a strongly summable, injective Dirichlet-convolution realization;
in particular, zeta has an exact Euler product and no zeros there. A
strongly summed Hurwitz expansion at an infinite positive shift recovers the
Stirling log-Gamma by differentiation. At infinite imaginary part, however,
phase choices matter. The published phase-forgetting surcomplex exponential,
coefficientwise Dirichlet summation on translated galaxies, and a single-pole
zeta continuation are incompatible. A corresponding functional-equation
obstruction and a Gamma reflection obstruction are proved.

An alternative, elementary-extension construction supports genuine
infinite-height zeros and transfers classical zero-free regions, counting
estimates and density results, but does not automatically agree with the
canonical surreal exponential or Hahn summation. Recent 2026 proportion
preprints are distinguished from published benchmarks. Finally, assuming RH
and simplicity, we prove that nonreal zeros created at negative infinitesimal
de Bruijn--Newman time must escape every ordinary bounded region.
\end{abstract}
\vfill
\noindent\small
\textbf{Status.} Complete arguments are supplied for the stated derived results,
relative to the explicitly cited classical and surreal foundations. This is an
unrefereed research manuscript, not a proof-assistant-verified development.
It does not assert a proof of the ordinary Riemann hypothesis, a unique global
surcomplex zeta function, or priority for every consequence derived here.
\end{titlepage}
\begingroup
\footnotesize
\setlength{\parskip}{0pt}
\tableofcontents
\endgroup
\newpage

\section{Scope, provenance, and the central distinction}
\label{sec:scope}

The question is not answered by replacing the symbol $\C$ with $\No[i]$ in a
classical formula. The formulas defining Gamma and zeta use several different
operations: ordinary limits, integrals, infinite products, analytic
continuation, complex exponentiation and arithmetic counting. These operations
need not have the same extensions to a non-Archimedean field. The choice of
extension is part of the definition, not a technicality that can be postponed.

We write $\SC=\No[i]$ for the surcomplex field. Four constructions will be kept
separate throughout:
\begin{center}
\begin{tabular}{P{0.20\textwidth}P{0.32\textwidth}P{0.36\textwidth}}
\toprule
Construction & Domain and mechanism & What it establishes\\
\midrule
Canonical finite lift & Finite arguments; ordinary germs evaluated at infinitesimals & Exact classical divisors and an RH equivalent to classical RH\\[3pt]
Infinite-right series & Positive infinite real part and finite imaginary part; strong Hahn sums & An injective arithmetic algebra, Euler products and zero-freeness\\[3pt]
Selected global prescription & Infinite-height sectors; a specified phase and continuation rule & Extra hypotheses are necessary; some natural combinations are inconsistent\\[3pt]
Elementary extension & A set-sized elementary model, embedded in $\No[i]$ & Transfer at infinite heights, with transported rather than automatically canonical operations\\
\bottomrule
\end{tabular}
\end{center}

\subsection{What was read in the repository}

The repository was inspected at the pinned snapshot
\begin{center}\small\pin.\end{center}
The principal inspected documents were the documentation catalogue and the
maintained guides to the Gamma-functions, surcomplex-analysis and
trigonometry reports \cite{repo,repoGamma,repoAnalysis,repoTrig}. The catalogue
explicitly distinguishes research drafts, finite verification programs and
partial Lean coverage. Its general warnings about strong summation,
coefficientwise summation, fine topology and coefficient rings are retained
here. This is a targeted audit, not an exhaustive review of every source file
or every theorem in the repository. A repository search returning no zeta
matches is not evidence that every existing zeta-related observation is absent.

In particular, the Gamma report already develops a Taylor--Stirling baseline,
proves nonuniqueness under literal surreal Bohr--Mollerup conditions, and
studies a finite-phase tube. Those results are not claimed as new here.
The finite-lift and phase foundations likewise have substantial antecedents
in the analysis and trigonometry reports and in the literature
\cite{CE,EK}. The present contribution is the Gamma--zeta comparison, the
explicit arithmetic and deformation constructions, and the exact logical
boundaries of possible RH generalizations.

\subsection{How to read the theorem claims}

The statements about ordinary Gamma, zeta, zero counting and the
de Bruijn--Newman constant are imported results. Hahn support lemmas and the
surreal exponential are also imported foundations. Results proved from these
inputs are called \emph{derived theorems}; this describes their proof status,
not a certified claim of historical novelty.

The most substantial derived statements are the robust RH criterion
(Theorem~\ref{thm:robust}), the arbitrary-coefficient Dirichlet realization
(Theorem~\ref{thm:Dirichlet}), the two global phase obstructions
(Theorems~\ref{thm:gamma-obstruction} and \ref{thm:resonance}), and the
conditional escape theorem for infinitesimal heat time
(Theorem~\ref{thm:escape}). Their proofs are independent of any assumed
numerical verification of RH.

\section{Finite parts, Hahn evaluation, and what a sum means}
\label{sec:foundations}

\subsection{Finite elements and infinitesimals}

The modulus on $\SC$ is $|x+iy|=(x^2+y^2)^{1/2}$. Define
\begin{align*}
\OO&=\{z\in\SC: |z|<r\text{ for some }r\in\R_{>0}\},\\
\mm&=\{z\in\SC: |z|<r\text{ for every }r\in\R_{>0}\}.
\end{align*}
Every $z\in\OO$ has a unique expression $z=c+\varepsilon$ with
$c\in\C$ and $\varepsilon\in\mm$. We write $c=\st z$. For an ordinary
set $U\subseteq\C$, its halo is $U^\#=\st^{-1}(U)$. Thus $\C^\#=\OO$.
The halo $c+\mm$ of one ordinary point is much smaller than a \emph{galaxy}
$b+\OO$; this distinction will matter at infinite height.

A real surreal has a normal-form decomposition
$x=\pp(x)+\fin(x)$, where $\pp(x)$ contains its purely infinite terms and
$\fin(x)$ is finite. This is not the same as standard part: $\fin(x)$ retains
infinitesimal terms. On finite real arguments there is a canonical trigonometric
map
\begin{equation}
\cis(r+\epsilon)=e^{ir}\Sh_{n\geq0}\frac{(i\epsilon)^n}{n!},
\qquad r\in\R,\quad \epsilon\in\mm\cap\No.
\label{eq:cis}
\end{equation}
It maps the additive group of finite real surreals onto the entire norm-one
surcomplex torus, with kernel $2\pi\Z$; see \cite{EK,repoTrig}.

\subsection{Support conventions}

In a set-sized Hahn workspace we write monomials as $t^\gamma$, with
larger $\gamma$ meaning smaller magnitude. One may regard
$t^\gamma=\omega^{-\gamma}$ when translating Conway normal forms. A Hahn
series has well-ordered exponent support. The natural valuation $v$ is its
least exponent; nonzero ordinary complex constants have valuation zero.
For nonzero $a,b$, the notation $a\prec b$ means $a/b\in\mm$, and
$a\sim b$ means $a/b\in1+\mm$.

A family is \emph{strongly summable} when the union of its exponent supports
is well ordered and each exponent occurs in only finitely many family
members. Its strong sum is denoted $\Sh$. Every support calculation in this
article is set-sized. Workspaces may be enlarged to contain the elements
and exponentials involved; no single fixed Hahn field is asserted to be
closed under every operation used here.

We use the usual Neumann support lemma: if $S$ is a well-ordered subset of the
positive cone of an ordered abelian group, the additive monoid generated by
$S$ is well ordered, and each group element has only finitely many
representations by finite multisets from $S$. This is the basic justification
for evaluating power series at infinitesimals. The standard Hahn and surreal
background is discussed in \cite{Conway,Gonshor,CE,repoAnalysis}.

\begin{lemma}[Infinitesimal formal evaluation]
\label{lem:evaluation}
Let $a_n\in\C$ and $\varepsilon\in\mm$. Then
$\Sh_{n\geq0}a_n\varepsilon^n$ exists. Evaluation commutes with sums,
products and compositions of formal power series when the inner series has
zero constant term. The corresponding assertion holds for finitely many
infinitesimal variables.
\end{lemma}
\begin{proof}
For $\varepsilon\ne0$, its exponent support is well ordered and positive.
Every exponent appearing in a power of $\varepsilon$ is a finite sum from
that support. Neumann's lemma gives both the support condition and local
finiteness of the contributions. Multiplication and composition can
therefore be regrouped exponent by exponent, with only finitely many terms
at each exponent. For several variables use the finite union of their
supports, which is again well ordered and positive. The case of a zero
variable is immediate.
\end{proof}

No growth bound on the ordinary coefficients $a_n$ is required. Their
Archimedean sizes do not change their Hahn valuations. Consequently a
classically divergent formal series can have an exact value at a nonzero
surreal infinitesimal. This is useful, but it does not turn that value into
an ordinary analytic limit.

\subsection{Three incompatible interpretations of an infinite series}

An ordinary convergent complex sum is computed in $\C$. A strong Hahn sum is
computed by support rules. A coefficientwise mixed sum first performs ordinary
complex sums at each coefficient and only then takes a strong Hahn sum. They
are different operations.

For example, at ordinary $\sigma>1$, the family $(n^{-\sigma})_{n\geq1}$
is not strongly summable: every member has exponent zero. Nevertheless its
ordinary complex sum is $\zeta(\sigma)$. Conversely, at positive infinite
$X$, the family $(n^{-X})_{n\geq1}$ has separated monomials and is strongly
summable. Neither assertion follows from the other.

A strong sum is not, merely by being a strong sum, the limit of finite partial
sums in the full surreal fine topology. Formal differentiation, ordinary
complex differentiation, fine differentiation of a local lift and the
Berarducci--Mantova scalar derivation are also distinct operations. We use
formal or local analytic differentiation below, and make no unstated scalar
chain-rule assertion for the Berarducci--Mantova derivation.

\section{Canonical Gamma, zeta, and xi on the finite plane}
\label{sec:finite}

\subsection{Taylor--Laurent extension}

\begin{definition}[Canonical finite lift]
For an ordinary holomorphic function $f$ near $c\in\C$, set
\begin{equation}
f^\#(c+\varepsilon)=\Sh_{n\geq0}
 \frac{f^{(n)}(c)}{n!}\varepsilon^n,
\qquad \varepsilon\in\mm.
\label{eq:lift}
\end{equation}
For an ordinary meromorphic germ
$f(c+u)=u^m\sum_{n\geq0}a_nu^n$, with $m\in\Z$ and $a_0\ne0$, use
\begin{equation}
f^\#(c+\varepsilon)=\varepsilon^m\Sh_{n\geq0}a_n\varepsilon^n
\label{eq:laurent}
\end{equation}
when $\varepsilon\ne0$. At $\varepsilon=0$ use the ordinary value if the
germ is holomorphic, and retain a meromorphic pole otherwise.
\end{definition}

Thus a pole at $c$ does not delete the entire halo $c+\mm$: every nonzero
infinitesimal displacement has a well-defined surcomplex value, usually of
infinite magnitude. A meromorphic pole is not itself a field element
called infinity.

\begin{proposition}[Functoriality of canonical lifting]
\label{prop:functorial}
The construction respects addition, multiplication, reciprocals where
nonzero, conjugation, differentiation, and composition where the underlying
ordinary germs can be composed. Every identity of ordinary meromorphic
functions lifts to an identity on its finite halo, with removable
singularities handled before evaluation.
\end{proposition}
\begin{proof}
These are identities of Taylor or Laurent series. Lemma~\ref{lem:evaluation}
permits their formal operations after infinitesimal evaluation. For a
reciprocal, the nonzero constant coefficient of a unit permits a geometric
series in its infinitesimal remainder. Composition is applied after
subtracting the ordinary value of the inner function. Differentiation is
coefficientwise formal differentiation of a local series; translating that
series by another infinitesimal uses the same support lemma. Conjugation
fixes the real monomials and conjugates ordinary coefficients.
\end{proof}

The word ``canonical'' in this section refers to the explicit Taylor--Laurent
rule, not to uniqueness among all functions agreeing with $f$ on $\C$.

\subsection{An exact divisor theorem}

\begin{theorem}[Divisor preservation and infinitesimal valuation]
\label{thm:divisor}
Let $f$ be a nonzero ordinary meromorphic function and let
$m=\ord_c f\in\Z$. Write
$f(c+u)=a_mu^m+\cdots$, $a_m\ne0$. For every nonzero
$\varepsilon\in\mm$,
\begin{equation}
f^\#(c+\varepsilon)\sim a_m\varepsilon^m,
\qquad v\bigl(f^\#(c+\varepsilon)\bigr)=m\,v(\varepsilon).
\label{eq:divisorvaluation}
\end{equation}
In particular, $f^\#$ has no zero or pole at a nonzero infinitesimal
displacement from $c$. Its zeros and poles on the finite plane are exactly
the ordinary zeros and poles of $f$, with the same multiplicities.
\end{theorem}
\begin{proof}
Factor the Laurent series as
$a_mu^m(1+u g(u))$, where $g\in\C[[u]]$. Its evaluation is
$a_m\varepsilon^m(1+\eta)$ with $\eta\in\mm$. The factor $1+\eta$
is a nonzero unit of valuation zero. This proves both formulas and excludes
zeros away from $\varepsilon=0$. The order at the centre is read directly
from the same local factorization. At any finite point the centre is uniquely
specified by standard part, so these local conclusions exhaust the domain.
\end{proof}

The theorem is stronger than preservation of standard parts of zeros.
It says that a fixed ordinary function does not acquire a cloud of displaced
zeros merely by adjoining infinitesimals. Displaced zero clusters require
an actual change of the function or its coefficients.

\subsection{Gamma and its identities}

The classical Gamma function is meromorphic, has no zeros, and has simple
poles at $0,-1,-2,\ldots$, with residues $(-1)^n/n!$ at $-n$;
its recurrence, reflection and multiplication identities are recorded in
\cite{DLMFgamma,DLMFidentities}. Hence:

\begin{corollary}
The canonical finite Gamma function $\Gp$ is zero-free. Its only poles are
the ordinary nonpositive integers, with their classical residues. It satisfies
\begin{align}
\Gp(z+1)&=z\Gp(z),\label{eq:gamma-rec}\\
\Gp(z)\Gp(1-z)&=\frac{\pi}{\sin^\#(\pi z)},\label{eq:gamma-ref}\\
\prod_{j=0}^{q-1}\Gp\!\left(z+\frac jq\right)
 &=(2\pi)^{(q-1)/2}q^{1/2-qz}\Gp(qz),
 \qquad q\in\N_{>0},\label{eq:gamma-gauss}
\end{align}
as identities of finite meromorphic lifts. The index $q$ is an ordinary
positive integer.
\end{corollary}

For example, at a nonzero infinitesimal $\varepsilon$,
\begin{equation}
\Gp(\varepsilon)=\varepsilon^{-1}-\gamma+
 \left(\frac{\gamma^2}{2}+\frac{\pi^2}{12}\right)\varepsilon
 +\bigO(\varepsilon^2),
\label{eq:gamma-zero}
\end{equation}
where $\gamma$ is Euler's constant. Here and below an $\bigO(\varepsilon^k)$
remainder in a displayed Taylor example means $\varepsilon^k$ times a
finite value obtained from the indicated ordinary germ; it is not an
unexplained asymptotic limit along a set-indexed surreal sequence.

\subsection{Zeta, its completed function, and finite functional equations}

Define $\zp$ by the same meromorphic rule. Its unique pole is the ordinary
point $1$, with residue $1$. At a nonzero infinitesimal displacement,
\begin{equation}
\zp(1+\varepsilon)=\varepsilon^{-1}
 +\Sh_{n\geq0}\frac{(-1)^n\gamma_n}{n!}\varepsilon^n,
\label{eq:zeta-pole}
\end{equation}
where $\gamma_n$ are the classical Stieltjes constants \cite{DLMFzeta}.
The trivial zeros are exactly the ordinary points $-2,-4,\ldots$, and are
simple. All other zeros are precisely the ordinary nontrivial zeros,
with their original multiplicities.

Use the standard normalization
\begin{equation}
\xi(s)=\frac12s(s-1)\pi^{-s/2}\Gamma(s/2)\zeta(s).
\label{eq:xi}
\end{equation}
It is entire, satisfies $\xi(s)=\xi(1-s)$ and
$\xi(\overline s)=\overline{\xi(s)}$, and has $\xi(0)=\xi(1)=1/2$;
see \cite{DLMFfe}. Define $\xip$ by lifting the entire function, not by
multiplying undefined pole values in \eqref{eq:xi}. The product formula
then holds meromorphically wherever its factors are interpreted, including
its removable extensions.

\begin{proposition}[Finite completed functional equation]
\label{prop:finitefe}
On the finite surcomplex plane,
\begin{equation}
\xip(s)=\xip(1-s),\qquad
\xip(\overline s)=\overline{\xip(s)}.
\end{equation}
The ordinary zeta functional equation also lifts:
\begin{equation}
\zp(s)=2^s\pi^{s-1}\sin^\#(\pi s/2)\Gp(1-s)\zp(1-s),
\label{eq:zetafe}
\end{equation}
as a meromorphic identity. All powers in this finite formula use the
canonical finite exponential.
\end{proposition}
\begin{proof}
Apply Proposition~\ref{prop:functorial} to the corresponding ordinary
identities. At points with apparent products of a pole and a zero, first
multiply the Laurent germs; their cancellation is a formal algebraic
operation with finitely many negative powers.
\end{proof}

\subsection{How the finite Dirichlet series must be interpreted}

Let $s=c+\varepsilon$ with $\Re c>1$. The correct mixed expression is
\begin{equation}
\zp(c+\varepsilon)=\Sh_{k\geq0}\frac{(-\varepsilon)^k}{k!}
 \left(\sum_{n\geq1}^{\C}(\log n)^k n^{-c}\right).
\label{eq:mixed}
\end{equation}
For each ordinary $k$, the inner series converges absolutely in $\C$ and
is the appropriate zeta derivative. Its value is an ordinary coefficient.
Only the outer series is a strong Hahn sum. Reading the unadorned expression
$\sum n^{-s}$ as a strong sum at finite $s$ would be incorrect: infinitely
many nonzero terms share exponent zero. This remains incorrect even when
$\Re s$ is a large ordinary real number.

\section{RH on finite halos and an infinitesimal stability criterion}
\label{sec:robust}

\subsection{The direct finite generalization}

\begin{definition}
The \emph{canonical finite surcomplex Riemann hypothesis}, denoted
$\RH_{\mathrm{fin}}$, is
\begin{equation}
\forall s\in\OO\quad
\bigl(\xip(s)=0\ \Longrightarrow\ \Re s=1/2\bigr).
\label{eq:RHfinite}
\end{equation}
\end{definition}

\begin{theorem}
\label{thm:finiteRH}
$\RH_{\mathrm{fin}}$ is equivalent to the ordinary Riemann hypothesis.
\end{theorem}
\begin{proof}
If the ordinary RH holds, every zero of $\xip$ is an ordinary zero of
$\xi$ by Theorem~\ref{thm:divisor}, and therefore has real part $1/2$.
Conversely, restrict \eqref{eq:RHfinite} to ordinary complex arguments.
\end{proof}

This equivalence is useful: it identifies the correct zero geometry and
prevents spurious ``infinitesimal counterexamples'' obtained by replacing an
exact zero with a merely infinitesimal function value. For example, if
$\rho$ is a simple ordinary zero and $\varepsilon\ne0$ is infinitesimal,
then $\xip(\rho+\varepsilon)\sim\xi'(\rho)\varepsilon$ is nonzero.
An infinitesimal residual is not a zero.

\subsection{Symmetry-preserving perturbations}

Let $\symh$ be the real vector space of ordinary entire functions $h$ with
\begin{equation}
h(1-s)=h(s),\qquad
h(\overline s)=\overline{h(s)}.
\label{eq:hsym}
\end{equation}
For real $\epsilon\in\mm$, set
\begin{equation}
F_{\epsilon,h}(s)=\xip(s)+\epsilon h^\#(s),\qquad s\in\OO.
\label{eq:perturb}
\end{equation}
This is a genuine coefficient perturbation of the function. It is not the
unperturbed lift studied by the divisor theorem.

\begin{theorem}[Exact critical-line stability at a simple zero]
\label{thm:simplestable}
Suppose $\rho$ is a simple zero of $\xi$ on $\Re s=1/2$ and $h\in\symh$.
For every real infinitesimal $\epsilon$, $F_{\epsilon,h}$ has exactly one
zero in $\rho+\mm$. That zero lies exactly on $\Re s=1/2$.
Its expansion begins
\begin{align}
s_{\epsilon}=\rho
 &-\epsilon\frac{h(\rho)}{\xi'(\rho)}\notag\\
 &+\epsilon^2\left(
 \frac{h(\rho)h'(\rho)}{\xi'(\rho)^2}
 -\frac{\xi''(\rho)h(\rho)^2}{2\xi'(\rho)^3}
 \right)+\bigO(\epsilon^3).
\label{eq:rootshift}
\end{align}
If $h(\rho)\ne0$ and $\epsilon\ne0$, the displacement has valuation
$v(\epsilon)$.
\end{theorem}
\begin{proof}
Consider the ordinary two-variable analytic germ
$A(s,u)=\xi(s)+u h(s)$ at $(\rho,0)$. Since
$\partial_sA(\rho,0)=\xi'(\rho)\ne0$, the ordinary implicit-function
theorem gives an analytic root germ $s=r(u)$ and an analytic factorization
\[
A(s,u)=(s-r(u))B(s,u),\qquad B(\rho,0)=\xi'(\rho)\ne0.
\]
Evaluate both infinitesimal variables by Lemma~\ref{lem:evaluation}.
The evaluated $B$ is a unit throughout the halo, proving existence and
uniqueness there.

The involution $J(s)=1-\overline s$ preserves the halo of $\rho$ and,
by \eqref{eq:hsym}, sends a zero of $F_{\epsilon,h}$ to a zero. Uniqueness
forces $J(s_\epsilon)=s_\epsilon$, which is the exact equality
$\Re s_\epsilon=1/2$. Substitution of
$r(u)=\rho+a_1u+a_2u^2+\cdots$ gives
$a_1=-h/\xi'$ and the stated $a_2$. If $h(\rho)\ne0$, $a_1$ is a
nonzero ordinary coefficient, so the valuation assertion follows.
\end{proof}

It is essential that the coefficient of the perturbation is real and that
both symmetries are retained. Without these conditions, even a simple
critical-line zero can move transversely.

\subsection{A characterization involving RH and simplicity}

\begin{theorem}[Robust finite RH criterion]
\label{thm:robust}
The following three statements are equivalent.
\begin{enumerate}[label=\textup{(\roman*)}]
\item The ordinary RH holds and every nontrivial zeta zero is simple.
\item For every ordinary real constant $c$ and every positive real
infinitesimal $\epsilon$, all finite zeros of $\xip+c\epsilon$ have
real part $1/2$.
\item For every $h\in\symh$ and every real infinitesimal $\epsilon$, all
finite zeros of $F_{\epsilon,h}$ have real part $1/2$.
\end{enumerate}
\end{theorem}
\begin{proof}
Assume (i). If $F_{\epsilon,h}(s)=0$ and $s$ is finite, taking standard
parts gives $\xi(\st s)=0$. By (i), this centre is a simple critical-line
zero. Theorem~\ref{thm:simplestable} gives $\Re s=1/2$. Thus (iii) holds,
and (iii) implies (ii) by taking $h$ constant.

Assume (ii). Taking $c=0$ and using Theorem~\ref{thm:finiteRH} proves RH.
Suppose, for a contradiction, that $\rho=1/2+i\gamma$ is a zero of
multiplicity $m\geq2$. The ordinary function
\[
 g(u)=\xi(\rho+iu)
\]
has real Taylor coefficients at $u=0$: it is real on ordinary real $u$,
by the two xi symmetries. Hence
\[
g(u)=b_m u^m+b_{m+1}u^{m+1}+\cdots,\qquad
b_m=\frac{i^m\xi^{(m)}(\rho)}{m!}\in\R\setminus\{0\}.
\]
Choose the constant $c=b_m$, choose a positive infinitesimal $\epsilon$,
and put $\delta=\epsilon^{1/m}>0$. Substitution
$s=\rho+i\delta v$ into $\xip(s)+b_m\epsilon=0$ and division by
$b_m\delta^m$ gives the analytic germ equation
\begin{equation}
v^m+1+\sum_{k>m}\frac{b_k}{b_m}\delta^{k-m}v^k=0.
\label{eq:split}
\end{equation}
Every root $r$ of $v^m+1$ is simple. The ordinary implicit-function theorem
at $(\delta,v)=(0,r)$, followed by Hahn evaluation, gives a solution
$v=r+\bigO(\delta)$ of \eqref{eq:split}. For $m\geq2$, at least one such
$r$ is nonreal. For that solution,
\[
\Re s=\frac12-\delta\Im v\ne\frac12,
\]
since $\Im v$ has nonzero ordinary standard part. This contradicts (ii).
Thus every zero is simple and (i) follows.
\end{proof}

\begin{corollary}[Multiplicity determines the splitting scale]
In the multiple-zero construction of the proof, the off-line displacement
has valuation $v(\epsilon)/m$. In particular, a symmetric perturbation of
order $\epsilon$ can produce a much larger transverse displacement when
$m>1$.
\end{corollary}

\begin{example}
At a hypothetical double zero, the leading local equation is
$u^2+\epsilon=0$. Its roots are $u=\pm i\sqrt\epsilon$, so
$s=\rho+iu$ moves horizontally by $\mp\sqrt\epsilon$. The perturbation
respects the functional-equation symmetries, but the two displaced zeros
are exchanged by reflection in the critical line instead of being individually
fixed by it.
\end{example}

The theorem packages two ordinary conjectural properties into one exact
non-Archimedean robustness statement. It does not prove either of those
ordinary properties. Its mechanism is the classical stability of simple
real zeros and the splitting of multiple zeros, made exact by Hahn
evaluation and the critical-line involution.

\subsection{Stable halos and proportions}

For a fixed symmetric perturbation, call an ordinary critical-line zero a
\emph{uniformly stable centre} when the conclusion of
Theorem~\ref{thm:simplestable} holds for every such $h$ and every real
infinitesimal parameter. Every simple critical-line zero is a uniformly
stable centre. The proof of Theorem~\ref{thm:robust} also shows that a
multiple critical-line zero is not one.

Consequently every lower bound on the classical proportion of simple
critical-line zeros is immediately a lower bound on the proportion of
uniformly stable centres. This is a count of centre-indexed halos, with the
classical total zero count as denominator. It is not, without a separate
boundary convention, a claim about the number of displaced roots below an
arbitrary surcomplex height. Section~\ref{sec:partial} supplies explicit
published and recent values.

\section{Dirichlet series at positive infinite real part}
\label{sec:dirichlet}

\subsection{A domain on which ordinary Dirichlet coefficients separate}

Let
\begin{equation}
\mathcal D_+=\{s\in\SC:\Re s>r\text{ for every }r\in\R,
\quad \Im s\text{ is finite}\}.
\label{eq:Dplus}
\end{equation}
On arguments with finite imaginary part there is an unambiguous
finite-phase exponential
\begin{equation}
E_{\mathrm{fp}}(x+iy)=e^x\cis(y),
\qquad x\in\No,\quad y\in\OO\cap\No,
\label{eq:finitephase}
\end{equation}
using the Gonshor real exponential. For $s\in\mathcal D_+$ define
$n^{-s}=E_{\mathrm{fp}}(-s\log n)$ for ordinary positive integers $n$.

Write
\begin{equation}
s=X+c+\varepsilon,
\quad X=\pp(\Re s)>0,
\quad c\in\C,
\quad\varepsilon\in\mm.
\label{eq:Ddecomp}
\end{equation}
Then
\begin{equation}
n^{-s}=q_n n^{-c}\Sh_{k\geq0}
 \frac{(-\varepsilon\log n)^k}{k!},
\qquad q_n=e^{-X\log n}.
\label{eq:qndecomp}
\end{equation}
The exponential of a purely infinite real surreal is a Conway monomial;
this standard property of the surreal exponential makes each $q_n$ a
monomial \cite{Gonshor,EK}. Moreover,
\[
q_mq_n=q_{mn},\qquad 1=q_1>q_2>q_3>\cdots,
\qquad q_n/q_m\in\mm\quad(n>m).
\]
The separation is stronger than ordinary convergence of a Dirichlet series.
Each successive ordinary index has a strictly smaller Archimedean scale.

\begin{theorem}[Arbitrary-coefficient Dirichlet realization]
\label{thm:Dirichlet}
For each fixed $s\in\mathcal D_+$ and every arithmetic function
$a:\N_{>0}\to\C$, the strong sum
\begin{equation}
\DD_s(a)=\Sh_{n\geq1}a(n)n^{-s}
\label{eq:Da}
\end{equation}
exists. No growth assumption on $a(n)$ is needed. The map $\DD_s$ is an
injective unital $\C$-algebra homomorphism from arithmetic functions with
Dirichlet convolution to $\SC$:
\begin{equation}
\DD_s(a*b)=\DD_s(a)\DD_s(b),
\qquad (a*b)(n)=\sum_{d\mid n}a(d)b(n/d).
\label{eq:conv}
\end{equation}
If $a\ne0$ and $n_0$ is its least nonzero index, then
\begin{equation}
\DD_s(a)\sim a(n_0)n_0^{-s}.
\label{eq:firstindex}
\end{equation}
\end{theorem}
\begin{proof}
The positive valuations $v(q_n)$, $n\geq2$, form a strictly increasing
sequence, hence a well-ordered set. Adjoin the positive, well-ordered
support of $\varepsilon$. Their finite union $S$ is well ordered and
positive. Formula~\eqref{eq:qndecomp} expresses every nonconstant
contribution using a finite sum of elements of $S$. Neumann's lemma
therefore proves strong summability, including the local finiteness of
contributions from different $n$ and different powers of $\varepsilon$.
Ordinary complex coefficients, however large as functions of $n$, do not
change these supports.

For multiplication, expand the two strongly summable families. The same
support argument permits regrouping by $mn=k$. For each ordinary $k$ there
are only finitely many ordered factorizations $mn=k$, and
$m^{-s}n^{-s}=k^{-s}$. This gives \eqref{eq:conv}. The arithmetic function
supported at $1$ with value $1$ maps to the field unit.

For the last assertion, the $n_0$ term has leading monomial $q_{n_0}$
with nonzero ordinary coefficient $a(n_0)n_0^{-c}$. Every term with
$n>n_0$ has strictly smaller leading monomial, and all its other monomials
are smaller still. Thus none can cancel that coefficient. Division by the
$n_0$ term gives $1$ plus an infinitesimal. In particular the image of any
nonzero $a$ is nonzero, proving injectivity.
\end{proof}

This embeds the entire algebra of ordinary complex arithmetic functions,
not merely a classical half-plane convergence class. The least nonzero
index is a leading-term invariant of its image. The image need not be the
whole field. The units of the arithmetic algebra are exactly the functions
with $a(1)\ne0$, whereas any nonzero image is invertible in the ambient
field even when its reciprocal is not in that image.

\subsection{Zeta, Euler products, and logarithmic differentiation}

\begin{definition}
On $\mathcal D_+$, define
\begin{equation}
\zD(s)=\Sh_{n\geq1}n^{-s}.
\end{equation}
The subscript records its exact strong-Dirichlet prescription.
\end{definition}

\begin{theorem}[Euler identities and zero-freeness]
\label{thm:euler}
For every $s\in\mathcal D_+$,
\begin{align}
\zD(s)&=\Ph_{p\ \mathrm{prime}}(1-p^{-s})^{-1},
\label{eq:Euler}\\
\zD(s)^{-1}&=\Sh_{n\geq1}\mu(n)n^{-s},
\label{eq:Mobius}\\
\log\zD(s)&=\Sh_{p\ \mathrm{prime}}\Sh_{k\geq1}\frac{p^{-ks}}{k},
\label{eq:logEuler}\\
-\frac{\zD'(s)}{\zD(s)}&=\Sh_{n\geq1}\Lambda(n)n^{-s}.
\label{eq:Mangoldt}
\end{align}
The logarithm in \eqref{eq:logEuler} is the infinitesimal logarithm at $1$.
The prime product means its strongly summable expansion. In particular,
\begin{equation}
\zD(s)\ne0,\qquad \zD(s)-1\sim2^{-s}.
\label{eq:Dzerofree}
\end{equation}
\end{theorem}
\begin{proof}
The formal Euler product expands into terms indexed by finite prime
factorizations. Unique factorization assigns exactly one such monomial to
each ordinary positive integer. The support proof of
Theorem~\ref{thm:Dirichlet} justifies this expansion and regrouping. The
Dirichlet identity $\mathbf1*\mu=\delta_1$ gives the reciprocal formula.
Here $\mathbf1(n)=1$ and $\delta_1$ is the convolution identity.

Since $\zD(s)-1$ is infinitesimal, its formal logarithm exists. Expanding
the logarithm of each Euler factor gives the prime-power sum, again with
strongly summable support. Equivalently the identity can be verified in the
formal Dirichlet algebra, coefficient by coefficient, and then evaluated.
Local differentiation is legitimate by the next proposition and gives
\eqref{eq:Mangoldt}; it also follows by applying $\DD_s$ to
$\log n=\sum_{d\mid n}\Lambda(d)$. Finally the first nonzero term of
$\zD-1$ is $2^{-s}$, and $\zD=1+$ an infinitesimal is nonzero.
\end{proof}

\begin{proposition}[Local analytic dependence]
\label{prop:Dderivative}
For an arithmetic function $a$ and an infinitesimal $h$,
\begin{equation}
\DD_{s+h}(a)=\Sh_{k\geq0}\frac{h^k}{k!}
 \left(\Sh_{n\geq1}a(n)(-\log n)^k n^{-s}\right).
\label{eq:Dtaylor}
\end{equation}
The double family is strongly summable. Hence the parenthesized expression
is the $k$th local analytic derivative of $\DD_s(a)$ for every ordinary
$k\geq0$.
\end{proposition}
\begin{proof}
Use \eqref{eq:qndecomp} and adjoin the positive support of $h$ to the
support set used in Theorem~\ref{thm:Dirichlet}. Expansion of
$e^{-h\log n}$ produces exactly \eqref{eq:Dtaylor}, with regrouping
permitted by Neumann's lemma. The Taylor coefficients characterize the
local analytic germ.
\end{proof}

At positive infinite real $X$, every term in
$(-1)^k\zD^{(k)}(X)$ is nonnegative, and at least one is positive.
Thus it is strictly positive for every ordinary $k\geq0$. This is an
exact local complete-monotonicity statement on the infinite-right real
domain, not a claim obtained by taking a classical limit $X\to\infty$.

\subsection{What right-hand asymptotics do not determine}

At positive infinite $X$,
\begin{equation}
e^{-X^2}\prec n^{-X}\qquad(n\geq2\text{ ordinary}).
\label{eq:Dirichletflat}
\end{equation}
Indeed $X^2-X\log n$ is positive infinite, so the ratio is infinitesimal.
Adding a term $e^{-X^2}$ is invisible after every fixed ordinary Dirichlet
scale. Consequently agreement to every fixed Dirichlet asymptotic order
is weaker than equality to the exact strong sum. This is the zeta analogue
of the distinction between a complete power asymptotic expansion and a
chosen Gamma normalization at infinity.

Theorem~\ref{thm:euler} is a genuine new domain result for the specified
extension. It does not address zeros in an infinite-height critical strip:
$\Re s$ is positive infinite throughout $\mathcal D_+$.

\section{Gamma at infinity and a Hurwitz--Stirling bridge}
\label{sec:gamma-infinite}

\subsection{The strong Stirling logarithm}

For a positive infinite real $A$, define
\begin{equation}
\LSt(A)=\left(A-\frac12\right)\log A-A+\frac12\log(2\pi)
 +\Sh_{k\geq1}\frac{B_{2k}}{2k(2k-1)}A^{1-2k},
\label{eq:Stirling}
\end{equation}
and $\Gamma_{\mathrm{St}}(A)=e^{\LSt(A)}$. We use the standard Bernoulli
convention
\begin{equation}
\frac{u}{e^u-1}=\sum_{n\geq0}B_n\frac{u^n}{n!},
\qquad B_1=-\frac12.
\label{eq:Bernoulli}
\end{equation}
The coefficients in \eqref{eq:Stirling} are the classical Stirling
coefficients \cite{DLMFstirling}. The strong interpretation at an infinite
argument is part of the surreal extension framework in \cite{CE,repoGamma}.
Here it follows immediately from Lemma~\ref{lem:evaluation} applied to
$A^{-1}$ that the displayed correction is well defined.

The same expression defines $\LSt(z)$ for $\Re z>0$ and $|z|$ infinite,
using the principal finite angle:
\begin{equation}
\Log z=\log|z|+i\Arg z,\qquad -\pi/2<\Arg z<\pi/2.
\label{eq:principalLog}
\end{equation}
The canonical finite-angle parametrization of the unit circle supplies
$\Arg z$ even when the coordinates of $z$ are infinite. Since $z^{-1}$ is
infinitesimal, the correction series remains strongly summable. Locally,
$\Log(z+h)-\Log z=\log(1+h/z)$ whenever $h$ is sufficiently small to
remain in the half-plane; this defines the corresponding analytic germ.

\begin{proposition}[Exact Stirling recurrence]
\label{prop:Strec}
On this infinite right-half-plane domain,
\begin{equation}
\LSt(z+1)-\LSt(z)=\Log z.
\label{eq:Strec}
\end{equation}
If $E$ is any additive-to-multiplicative surcomplex exponential extending
$E_{\mathrm{fp}}$, then
$\Gamma_{E,\mathrm{St}}(z)=E(\LSt(z))$ is nonzero and satisfies
\begin{equation}
\Gamma_{E,\mathrm{St}}(z+1)=z\Gamma_{E,\mathrm{St}}(z).
\end{equation}
\end{proposition}
\begin{proof}
The Bernoulli identity proof of the Hurwitz shift relation in
Theorem~\ref{thm:Hurwitz}, differentiated at its parameter $s=0$, gives
\eqref{eq:Strec} as a formal identity in $z^{-1}$ and $\Log z$. Its proof
is formal and therefore applies also to complex $z^{-1}$, with the local
branch identity just noted. All series are strongly evaluable.
Exponentiation uses the homomorphism property and
$E(\Log z)=z$, since $\Arg z$ is finite. A homomorphism into
$\SC^\times$ never takes the value zero.
\end{proof}

\subsection{Why a phase condition remains necessary}

The logarithm $\LSt(z)$ need not have finite imaginary part. Thus the
existence of a canonical principal logarithm does not by itself make
$E_{\mathrm{fp}}(\LSt(z))$ meaningful. On the useful tube
$z=x+iy$ with $x>0$ infinite and $y$ finite, direct expansion gives
\begin{equation}
\Im\LSt(x+iy)=y\log x+R,
\qquad R\prec y\quad(y\ne0).
\label{eq:phaseleading}
\end{equation}
For example, the first correction is $-y/(2x)$; subsequent terms have still
smaller relative magnitude. Hence
\begin{equation}
\Im\LSt(x+iy)\text{ is finite}
\quad\Longleftrightarrow\quad y\log x\text{ is finite}.
\label{eq:phasetube}
\end{equation}
For $y=0$ the condition is automatic. This is the finite-phase tube phenomenon
already developed in the repository Gamma report \cite{repoGamma}.

For completeness, \eqref{eq:phaseleading} follows by writing
$\Log(x+iy)=\log x+\log(1+iy/x)$. The terms $y$ from
$(x-1/2)\Im\log(1+iy/x)$ and $-iy$ cancel, leaving
$y\log x-y/(2x)$ and terms divisible by $y/x^2$ or $y^3/x^2$.
The differentiated Stirling corrections are smaller still. Because $y$
is finite and $x$ is infinite, division of $R$ by a nonzero $y$ gives an
infinitesimal. An infinite $y\log x$ cannot be cancelled by such an $R$.

Outside \eqref{eq:phasetube}, a global exponential is available after
selecting its phase rule, but that selection has mathematical consequences.
Section~\ref{sec:phases} discusses the published canonical phase and the
obstructions it creates for a naive global zeta theory.

\subsection{The infinite-shift Hurwitz function}

The Hurwitz zeta function provides a direct bridge between zeta-type sums
and log-Gamma. Its classical identities and asymptotic expansion are
recorded in \cite{DLMFhurwitz}. We now interpret that expansion as an exact
strong sum at an infinite positive shift, rather than as an ordinary
remainder statement.

For positive infinite real $A$ and ordinary real $s\ne1$, define
\begin{equation}
Z_\infty(s,A)=\frac{A^{1-s}}{s-1}+\frac12 A^{-s}
 +\Sh_{k\geq1}\frac{B_{2k}}{(2k)!}
 (s)_{2k-1} A^{-s-2k+1},
\label{eq:Hurwitz}
\end{equation}
where $(s)_j=s(s+1)\cdots(s+j-1)$ and $(s)_0=1$.
For fixed ordinary $s$, the tail is an ordinary-coefficient power series
in $A^{-1}$ times $A^{-s}$ and is strongly summable.

A parameter germ at any ordinary real $c$ is also defined: put $s=c+h$,
require $h\log A$ to be infinitesimal, and use
\[
A^{-c-h}=A^{-c}\Sh_{j\geq0}\frac{(-h\log A)^j}{j!}.
\]
At $c=1$ this is a meromorphic parameter germ with a simple pole.
This restricted local parameter domain is sufficient for differentiation
at $s=0$. We do not assert an unrestricted finite-complex-parameter
prescription when its infinite phase has not been specified.

\begin{theorem}[Hurwitz shift, Bernoulli values, and log-Gamma]
\label{thm:Hurwitz}
For the construction \eqref{eq:Hurwitz}, the following identities hold
as meromorphic parameter identities, with removable products evaluated by
cancellation and parameter derivatives taken in the just-defined local germs:
\begin{align}
Z_\infty(s,A)-Z_\infty(s,A+1)&=A^{-s},
\label{eq:Hshift}\\
\partial_A Z_\infty(s,A)&=-s Z_\infty(s+1,A),
\label{eq:HderA}\\
Z_\infty(-m,A)&=-\frac{B_{m+1}(A)}{m+1},
\qquad m\in\N,\label{eq:Hnegative}\\
\left.\partial_s Z_\infty(s,A)\right|_{s=0}
 &=\LSt(A)-\frac12\log(2\pi).
\label{eq:Hgamma}
\end{align}
The residue at the parameter pole $s=1$ is $1$.
\end{theorem}
\begin{proof}
Put $D=\partial_A$. On the formal generalized-power expressions involved,
choose $D^{-1}A^{-s}=A^{1-s}/(1-s)$, with no added integration constant.
The Bernoulli generating identity gives
\begin{equation}
(1-e^D)^{-1}=-D^{-1}+\frac12
 -\sum_{k\geq1}\frac{B_{2k}}{(2k)!}D^{2k-1}.
\label{eq:Bernoperator}
\end{equation}
This is a formal Laurent identity in the symbol $D$; multiplication by
$1-e^D$ is interpreted before application. Since
\[
D^{2k-1}A^{-s}=-(s)_{2k-1}A^{-s-2k+1},
\]
its right-hand side applied to $A^{-s}$ is exactly \eqref{eq:Hurwitz}.
Also $e^D$ acts as the formal shift $A\mapsto A+1$, by the binomial
series in $A^{-1}$. Multiplying \eqref{eq:Bernoperator} by $1-e^D$
therefore proves \eqref{eq:Hshift}. Strong evaluation preserves this
formal identity. The argument is meromorphic in $s$, so covers the
parameter germs as well.

Termwise differentiation proves \eqref{eq:HderA}: the leading term
differentiates to $-A^{-s}$, the next to $-sA^{-s-1}/2$, and in the tail
$(s)_{2k-1}(s+2k-1)=s(s+1)_{2k-1}$.

For $s=-m$, the rising factorial vanishes as soon as it contains the
factor $s+m=0$, so the tail becomes finite. Expand
\[
B_{m+1}(A)=\sum_{j=0}^{m+1}\binom{m+1}{j}B_jA^{m+1-j}.
\]
Its $j=0$ and $j=1$ terms yield the first two terms of
\eqref{eq:Hurwitz}; the odd Bernoulli numbers beyond $B_1$ vanish.
For $j=2k\leq m+1$, the coefficient agrees because
$(-m)_{2k-1}=-m!/(m-2k+1)!$. This proves \eqref{eq:Hnegative}.

At $s=0$, differentiation of the first two terms gives
$A(\log A-1)-\tfrac12\log A$. Moreover,
\[
\left.\frac{d}{ds}(s)_{2k-1}\right|_{s=0}=(2k-2)!,
\]
so each tail coefficient becomes
$B_{2k}/[2k(2k-1)]$. This is precisely
\eqref{eq:Hgamma}. Differentiating \eqref{eq:Hshift} at zero now gives
\eqref{eq:Strec}, as used in Proposition~\ref{prop:Strec}.
Finally only the first term in \eqref{eq:Hurwitz} has a pole at $s=1$,
and its numerator is $1$ there.
\end{proof}

\begin{example}
The first special values are exact polynomials:
\[
Z_\infty(0,A)=\frac12-A,\qquad
Z_\infty(-1,A)=-\frac12A^2+\frac12A-\frac1{12}.
\]
They are not approximate large-$A$ formulas. The termination of the
rising factorial makes them exact for the specified strong construction.
\end{example}

\begin{warning}[The ordinary shifted sum is not the strong construction]
For ordinary real $s>1$ and positive infinite $A$, the family
$((A+n)^{-s})_{n\geq0}$ is not strongly summable. Indeed
$(A+n)^{-s}=A^{-s}(1+n/A)^{-s}$, and every summand has the same nonzero
leading term as $A^{-s}$. Thus \eqref{eq:Hurwitz} is not justified by
writing a strong sum of those terms. It is an exact asymptotic-series
prescription with the shift identities proved separately. Nor may one
substitute $A=1$ into an infinite-shift definition and thereby claim a
global Riemann zeta function.
\end{warning}

\subsection{Nonuniqueness is an existing Gamma phenomenon}

On positive finite arguments use the canonical Gamma lift, and on positive
infinite arguments use $\Gamma_{\mathrm{St}}$. The resulting baseline can
be modified on infinite galaxies without changing its finite values.
For example, if $M(x)$ is the leading infinite Conway monomial of $x$
with its real coefficient stripped, put
\begin{equation}
h_\lambda(x)=\lambda e^{-M(x)}\pp(x)
\quad(x>0\text{ infinite}),\qquad
h_\lambda(x)=0\quad(x>0\text{ finite}).
\label{eq:gauge}
\end{equation}
Multiplying the baseline by $e^{h_\lambda(x)}$ preserves the ordinary
unit-shift recurrence: $M(x+1)=M(x)$ and $\pp(x+1)=\pp(x)$.
Also $h_\lambda(qx)=q h_\lambda(x)$ and
$h_\lambda(x+j/q)=h_\lambda(x)$ for ordinary integers $q>0,j$.
Therefore any finite Gauss identities of the baseline are preserved.

The stronger assertion that this family can retain global strict
log-convexity, and the sharp classification for its larger monomialwise
family, belong to the repository Gamma report \cite{repoGamma}; they are
not reproved here. Its example already rules out a naive uniqueness claim
based on the classical Bohr--Mollerup axioms. The logarithmic correction
\eqref{eq:gauge} is smaller than every fixed inverse power at positive
infinite $x$, yet need not produce a small \emph{absolute} difference of
Gamma values. Full power asymptotics, relative error and absolute error
must be distinguished.

\section{Infinite phases and two obstruction theorems}
\label{sec:phases}

\subsection{A global exponential exists, but its phase is extra structure}

It would be incorrect to claim that there is no published global surcomplex
exponential. Ehrlich and Kaplan construct one, with a distinguished choice
associated to the omnific integer part \cite[Section 11]{EK}. In the
normal-form convention used here it is
\begin{equation}
E_0(x+iy)=e^x\cis(\fin y),
\qquad x,y\in\No.
\label{eq:E0}
\end{equation}
It extends \eqref{eq:finitephase}, is an additive-to-multiplicative
homomorphism onto $\SC^\times$, and has kernel $2\pi i\Oz$. We call its
phase \emph{phase-forgetting}: for every purely infinite real $T$ and
every ordinary real $r$,
\begin{equation}
E_0(irT)=1.
\label{eq:forget}
\end{equation}
This normalization is canonical relative to the chosen surreal integer-part
structure, but agreement with ordinary exponentiation does not uniquely
force it. Other phase extensions exist; the repository treats their
classification and infinite periods \cite{repoTrig}.

To separate the relevant datum, write $U(y)=E(iy)$ for a chosen homomorphism
from the additive real surreal group to the norm-one torus, extending finite
$\cis$. Since finite $\cis$ is onto that torus, every infinite positive
$T$ admits a finite $\theta$ with $U(T)=\cis\theta$. Thus
$P=T-\theta$ is a positive infinite period of $U$. The existence of such
periods is already a known structural fact; our first obstruction is its
application to Gamma's divisor.

\begin{theorem}[Gamma reflection forces extra infinite poles]
\label{thm:gamma-obstruction}
Let $U$ be a phase homomorphism extending finite $\cis$, and put
\[
\operatorname{Sin}_U(y)=\frac{U(y)-U(-y)}{2i}\qquad(y\in\No).
\]
There is no function $G$ taking a value in $\SC$ at every argument except
possibly the \emph{ordinary} nonpositive integers and satisfying
\begin{equation}
G(x)G(1-x)\operatorname{Sin}_U(\pi x)=\pi
\label{eq:globalreflection}
\end{equation}
whenever both displayed $G$ values are defined.
\end{theorem}
\begin{proof}
Choose a positive infinite period $P$ of $U$ as above, and set $x=P/\pi$.
Both $x$ and $1-x$ are infinite, so neither is an ordinary nonpositive
integer. Both $G$ values are therefore field elements. But
$\operatorname{Sin}_U(\pi x)=0$, making the left-hand side of
\eqref{eq:globalreflection} zero, a contradiction.
\end{proof}

The conclusion is not that global Gamma functions are impossible. A global
reflection formula must instead permit additional infinite poles, restrict
its domain, or change some other part of the prescription. Elementary
extensions do add nonstandard integer poles, as discussed below. The
ordinary nonpositive integers are the complete pole set only for the
\emph{finite canonical lift}.

\subsection{A simultaneous-resonance test}

\begin{proposition}[A two-frequency first-order distinction]
\label{prop:twofreq}
The following statement is true for ordinary real $t$ and remains true in
any elementary extension of the real structure with its ordinary cosine and
sine:
\begin{equation}
\cis(t\log2)=1\ \text{ and }\ \cis(t\log3)=1
\quad\Longrightarrow\quad t=0.
\label{eq:twofreq}
\end{equation}
It fails for the phase $U_0(t)=E_0(it)$ on $\No$.
\end{proposition}
\begin{proof}
Over $\R$, the two equalities imply
$t\log2=2\pi k$ and $t\log3=2\pi l$ for integers $k,l$.
Thus $k\log3=l\log2$, or $3^k=2^l$. Unique prime factorization, after
clearing negative exponents if needed, gives $k=l=0$ and $t=0$.
The implication is a first-order statement in the named real functions and
constants, so transfers to every elementary extension. In contrast, every
nonzero purely infinite $T$ satisfies both antecedents for $U_0$, by
\eqref{eq:forget}.
\end{proof}

Thus the phase-forgetting exponential cannot simply be treated as the full
complex exponential of an elementary extension. In particular, the presence
of a well-behaved real surreal exponential does not justify transferring
statements involving arbitrary infinite imaginary phases.

\subsection{Resonant galaxies and the zeta functional equation}

Fix a positive purely infinite $T$. On the translated finite galaxy
$iT+\OO$, write $s=iT+w$ with $w$ finite. For every ordinary positive
integer $n$, \eqref{eq:forget} gives
\begin{equation}
E_0(-s\log n)=E_0(-w\log n)=n^{-w}.
\label{eq:resonantn}
\end{equation}
Therefore coefficientwise ordinary Dirichlet summation on
$\Re\st w>1$ forces the value $\zp(w)$, not a new high-frequency function.

For clarity, a \emph{coherent meromorphic continuation in the galaxy
coordinate} means a section with one set-sized well-ordered Hahn support
and ordinary meromorphic coefficient functions on the connected ordinary
$w$-plane. Evaluation at finite $w$ uses Taylor--Laurent lifting. This is
the common-domain notion employed in the repository; it is stronger than
being fine-analytic separately at each point.

\begin{theorem}[Phase-forgetting zeta obstruction]
\label{thm:resonance}
Suppose $Z$ uses $E_0$ for its Dirichlet powers, agrees with coefficientwise
ordinary Dirichlet summation on the subdomains
$\Re\st w>1$ of both galaxies $iT+\OO$ and $-iT+\OO$, and has coherent
meromorphic continuations on their connected ordinary coordinate planes.
Then
\begin{equation}
Z(iT+w)=\zp(w),\qquad Z(-iT+w)=\zp(w)
\label{eq:replica}
\end{equation}
as meromorphic functions of finite $w$. In particular:
\begin{enumerate}[label=\textup{(\roman*)}]
\item $Z$ has poles at $1+iT$ and $1-iT$. It cannot have only the pole $1$.
\item $Z(3+iT)=\zeta(3)\ne0$, whereas $Z(-2-iT)=0$.
\item The usual pointwise zeta functional equation using $E_0$ cannot
hold at $s=3+iT$ with a Gamma factor that is a field-valued function at
$1-s=-2-iT$.
\end{enumerate}
\end{theorem}
\begin{proof}
Equation~\eqref{eq:resonantn} and the specified mixed summation give
$Z(iT+w)=\zp(w)$ in the ordinary right half-plane of $w$.
Subtract the canonical section of $\zeta$ from the coherent continuation.
For every ordinary $w$ in that half-plane, each Hahn coefficient of the
difference vanishes. The ordinary meromorphic identity theorem makes each
coefficient zero on the connected ordinary plane. This proves the first
identity in \eqref{eq:replica}; the second follows identically.

The pole and value assertions are now ordinary zeta facts. At $s=3+iT$,
the proposed equation would read
\[
\zeta(3)=2^s\pi^{s-1}\operatorname{Sin}_{E_0}(\pi s/2)
 G(-2-iT)\,Z(-2-iT)=0,
\]
where $\operatorname{Sin}_{E_0}(z)=(E_0(iz)-E_0(-iz))/(2i)$ and the
powers also use $E_0$. Every factor is a field element by hypothesis,
while the last factor is zero. This is impossible.
\end{proof}

The Gamma qualification is substantial, not vacuous. The shifted Stirling
prescription is field-valued and nonzero at the test point: define
\begin{equation}
G(-2-iT)=
\frac{\Gamma_{E_0,\mathrm{St}}(1-iT)}{(-2-iT)(-1-iT)(-iT)}.
\label{eq:shiftedGamma}
\end{equation}
The numerator is defined in the infinite right half-plane, and the three
denominator factors are nonzero. Thus using this natural Gamma extension
does not repair the contradiction. Introducing an actual Gamma pole there
would change the pole prescription and require a meromorphic cancellation,
not a pointwise multiplication of field values.

\begin{remark}[Exact scope of the obstruction]
The theorem rules out a conjunction of explicit axioms. It does not exclude
all global surcomplex zeta functions. A different phase, a different
summation rule, or a failure of connected ordinary-coordinate coherence
can avoid its assumptions. In fact, any phase $U$ having a common
nonzero infinite resonance $U(T\log n)=1$ for all ordinary $n$ encounters
the same argument. The existence of an infinite period of $U$ alone does
\emph{not} imply such a simultaneous logarithmic resonance.
\end{remark}

This also explains why the phrase ``trivial zeros'' needs care in a global
theory. Extra zeros of a chosen global sine are not automatically the
trivial-zero divisor of a completed zeta function. Gamma poles and zeta
zeros must cancel in a specified meromorphic construction before its
completed function and RH can be formulated.

\section{What weak global analyticity does not decide}
\label{sec:weakglobal}

The finite canonical lift fixes values on $\OO$ but leaves every infinite
galaxy untouched. Ordinary complex connectedness does not connect these
galaxies. Even exact functional-equation symmetries do not determine the
missing zeros.

\begin{theorem}[Explicit symmetric completions with different infinite zeros]
\label{thm:weakcompletions}
There are two class functions on $\SC$, locally Hahn-analytic at every
point, that agree with $\xip$ on the finite plane and satisfy
\begin{equation}
F(1-s)=F(s),\qquad F(\overline s)=\overline{F(s)},
\label{eq:weaksym}
\end{equation}
but have different infinite zero sets. One has no infinite zeros. The other
has exactly four infinite zeros, all simple and off the critical line.
\end{theorem}
\begin{proof}
Set
\[
F_0(s)=
\begin{cases}
\xip(s),&s\in\OO,\\
1,&s\notin\OO.
\end{cases}
\]
Finite translation and conjugation preserve finiteness, so the symmetries
hold. Every point has a sufficiently small fine neighborhood contained in
its galaxy. On the finite galaxy use the canonical analytic lift; elsewhere
$F_0$ is constant. This proves local Hahn analyticity and the absence of
infinite zeros.

Choose a positive purely infinite $T$ and an ordinary real $a$ with
$0<a<1/2$. Let $c_\pm=1/2\pm iT$. Replace $F_0$ on the two galaxies
$\pm iT+\OO$ by
\[
F_1(s)=(s-c_\pm)^2-a^2
\quad\text{on }\ \pm iT+\OO,
\]
and keep $F_1=F_0$ elsewhere. Conjugation exchanges the two polynomials.
Under $s\mapsto1-s$, the centred coordinate changes sign and the two
galaxies are exchanged, so evenness of the square gives \eqref{eq:weaksym}.
Local analyticity is unchanged. The added zeros are exactly
$c_\pm+a$ and $c_\pm-a$, all simple; their real parts are $1/2\pm a$.
\end{proof}

The two functions are not proposed as equally good zeta completions. In
particular no Euler product or Gamma factorization is asserted for them.
They prove the narrower, important point that finite agreement, local
analyticity and the two completed functional-equation symmetries do not
settle the infinite-height RH question. A global RH must name its global
function and the extra axioms that select it.

\section{An elementary-extension route to infinite heights}
\label{sec:transfer}

\subsection{The analytic structure that is transferred}

Start with the ordinary real field expanded by specified real functions.
To avoid partial-function ambiguities, one can name the real and imaginary
parts of the entire functions $\xi(s)$, $(s-1)\zeta(s)$ and
$1/\Gamma(s)$, treating a complex input as a pair of real coordinates.
Name the real exponential, ordinary sine and cosine, and whatever counting
functions or arithmetic predicates are needed. A meromorphic function is
then recovered by division off the zero set of its denominator. This is a
set-sized first-order language and a genuine ordinary structure.

Take a non-Archimedean elementary extension $\mathcal M^*$, for example a
nonprincipal ultrapower. Its underlying ordered field $F$ is real closed
and contains $\R$. It carries transferred functions $\xi^*$, $\zeta^*$,
$\Gamma^*$ and so forth. These are the operations to which transfer applies.

\begin{proposition}[An ordered-field realization inside the surreals]
\label{prop:embedding}
Every set-sized real closed ordered field extension $F$ of $\R$ admits an
ordered-field embedding into $\No$ fixing $\R$. Hence its complexification
embeds into $\SC$, and all the named transferred functions can be transported
to that embedded subfield.
\end{proposition}
\begin{proof}
Well-order a transcendence basis of $F$ over $\R$. Suppose a real closed
subfield $E$ has already been embedded, and let $\alpha$ be the next
transcendental basis element. The two sets
\[
L=\{j(e):e\in E,\ e<\alpha\},\qquad
R=\{j(e):e\in E,\ e>\alpha\}
\]
form a set-sized surreal cut. A surreal realizing $\{L\mid R\}$ lies
strictly between the two sets. It cannot belong to $j(E)$, since that would
put it on one of the two sides of its own cut. As $j(E)$ is real closed,
it is therefore transcendental over $j(E)$.

Sending $\alpha$ to this cut realization extends the embedding to an
ordered embedding of rational-function fields: signs of polynomials over
a real closed field are determined by the cut relative to their real
roots and the signs of their irreducible quadratic factors. Extend further
to the compatible real closure inside $\No$. At limit stages take unions.
After the whole basis is processed, real closedness and algebraicity of
$F$ over the resulting rational-function field give the desired embedding.
Adjoin $i$ to complexify. Transport of the extra named functions is then
by conjugation with this embedding.
\end{proof}

This construction supplies actual surcomplex-valued models with
infinite-height arguments. It does not supply a unique function on all of
$\SC$. Moreover an ordered-field embedding alone need not preserve the
Gonshor exponential, strong Hahn sums, or the canonical $E_0$ phase.
Proposition~\ref{prop:twofreq} gives a direct obstruction to identifying
the transferred full phase with $E_0$ on a domain containing a purely
infinite simultaneous resonance. These operations must not be silently
interchanged after embedding.

\subsection{RH transfers exactly, not heuristically}

\begin{theorem}[Elementary RH equivalence]
\label{thm:transferRH}
For any elementary extension of the specified xi structure,
\begin{equation}
\forall s\in F[i]\quad
\bigl(\xi^*(s)=0\Longrightarrow\Re s=1/2\bigr)
\label{eq:RHstar}
\end{equation}
is equivalent to the ordinary RH. The extension has zeros at genuinely
infinite heights. Every finite zero of $\xi^*$ is nevertheless an exact
ordinary zero, not merely a point infinitesimally close to one.
\end{theorem}
\begin{proof}
The assertion \eqref{eq:RHstar} is the transfer of a first-order sentence
about two real coordinates and the named components of $\xi$. Elementarity
preserves its truth in both directions.

Ordinary xi zeros have unbounded positive imaginary part. The first-order
statement that for every real $T$ there is a zero of imaginary part greater
than $T$ transfers to every $T\in F$, including infinite $T$. Thus the
extension contains infinite-height zeros.

Finally, fix an ordinary $R>0$. The ordinary closed disk $|s|\leq R$
contains a finite list of xi zeros $\rho_1,\ldots,\rho_m$. The assertion
that every zero in this disk equals one of those finitely many named points
is first order and true in the ordinary structure, hence remains true in
$\mathcal M^*$. Every finite $s\in F[i]$ belongs to some such ordinary
disk. This proves the last assertion.
\end{proof}

This finite-zero agreement does not assert equality of all finite values
$\xi^*(c+\epsilon)=\xip(c+\epsilon)$ after an arbitrary field embedding.
Transferred Taylor remainder estimates give all fixed finite-order
approximations. Exact equality with a specified strong Hahn sum is an
additional summation compatibility requirement. A value smaller than every
fixed power of a particular infinitesimal is not forced to be zero in a
non-Archimedean field.

If the integer predicate is included, the transferred reciprocal Gamma
has zeros at the transferred nonpositive integers, including nonstandard
ones. Thus $\Gamma^*$ acquires infinite poles. This is precisely one of
the alternatives allowed by Theorem~\ref{thm:gamma-obstruction}; the
reflection formula is not being combined with an incorrectly restricted
ordinary pole set.

\subsection{Counting functions and standard quantifiers}

Let $N(T)$ count ordinary nontrivial zeta zeros with $0<\Im\rho\leq T$,
with multiplicity. Naming $N$ in the language yields a transferred function
$N^*(T)$. At infinite $T$ its value is a nonstandard integer. It is not
the external cardinality of the set of ordinary zeros below an infinite
surreal height; that external set is already all ordinary positive-height
zeros and is countably infinite.

A classical estimate $A(T)=\bigO(B(T))$ must be rewritten as a quantified
inequality with witnesses: there are ordinary $C,T_0>0$ such that
$|A(T)|\leq C B(T)$ for all $T\geq T_0$. Keeping these ordinary witnesses,
transfer yields the same inequality for infinite $T$ in the model.
Likewise a bound for each fixed ordinary $\sigma$ and each fixed ordinary
$\eta>0$ does not automatically permit $\sigma$ or $\eta$ to be
infinitesimal perturbations of endpoints. Uniformity must occur in the
original theorem before it is transferred.

\section{How classical partial results generalize}
\label{sec:partial}

This section gives selected precise benchmarks and their exact scopes.
It is not a claim that every quoted numerical constant is the best currently
available one. The literature check includes August and September 2026
proportion preprints; these are marked separately from published results.

\subsection{Zero-free regions and numerical verification}

\begin{proposition}[Halo transfer of zero-free sets]
\label{prop:halofree}
If an ordinary set $A\subseteq\C$ contains no zeros of an ordinary
meromorphic function $f$, then $f^\#$ has no zeros on $A^\#$. Poles are
understood meromorphically and are not asserted to be function values.
\end{proposition}
\begin{proof}
Any zero of $f^\#$ is an exact ordinary zero of $f$, by
Theorem~\ref{thm:divisor}; its standard part is itself.
\end{proof}

Thus the classical zero-free half-plane $\Re s\geq1$, the critical-strip
localization of nontrivial zeros, and every verified finite-height zero-free
region hold exactly for the canonical finite lift \cite{DLMFzeros,Titch}.
There is no additional infinitesimal boundary layer of new zeros.

Platt and Trudgian rigorously verified RH for all zeros with
$0<\Im\rho\leq H$, where $H=3\cdot10^{12}$, using interval arithmetic
\cite{PT}. The same statement holds for $\xip$ at every finite surcomplex
argument in that height range. This is a cited verified benchmark, not a
new computation performed here and not a claim about the latest numerical
record. It gives no information about infinite heights. We do not infer
simplicity merely from the quoted verification statement.

For an explicit analytic benchmark, Mossinghoff, Trudgian and Yang prove
zero-freeness in
\begin{equation}
\sigma\geq1-
\frac{1}{55.241(\log|t|)^{2/3}(\log\log|t|)^{1/3}},
\qquad |t|\geq3,
\label{eq:MTY}
\end{equation}
and also in
\begin{equation}
\sigma\geq1-\frac{1}{5.558691\log|t|},\qquad |t|\geq2
\label{eq:MTYclassical}
\end{equation}
\cite{MTY}. Proposition~\ref{prop:halofree} gives the halo versions.
In an elementary extension, the actual inequalities
\eqref{eq:MTY}--\eqref{eq:MTYclassical} transfer to infinite $t$ as well,
using the \emph{transported} logarithm. Their widths are then infinitesimal
positive quantities. Substituting the Gonshor logarithm after an arbitrary
ordered-field embedding would require compatibility not supplied by the
embedding theorem.

At infinitely large positive real part the stronger direct result
\eqref{eq:Dzerofree} is already available for $\zD$, without transfer.
This is a different sector from \eqref{eq:MTY} at infinite height and
bounded real part.

\subsection{Counting zeros at infinite height}

The classical Riemann--von Mangoldt estimate is
\begin{equation}
N(T)=\frac{T}{2\pi}\log\frac{T}{2\pi}-\frac{T}{2\pi}
 +\bigO(\log T),\qquad T\geq2,
\label{eq:RvM}
\end{equation}
with multiplicities and a fixed endpoint convention \cite{Titch,DLMFzeros}.
Changing the convention at a zero height is absorbed by the displayed
error. The canonical finite lift has exactly this ordinary zero count at
every ordinary height.

When $N$ is named in the elementary structure, there is an ordinary
constant $C$ such that, at all sufficiently large transferred heights,
\begin{equation}
\left|N^*(T)-\frac{T}{2\pi}\log^*\frac{T}{2\pi}
             +\frac{T}{2\pi}\right|
 \leq C\log^*T.
\label{eq:RvMstar}
\end{equation}
In particular this holds for every infinite $T$ in the model. The normalized
ratio of $N^*(T)$ to its leading expression differs from $1$ by an
infinitesimal. Equation~\eqref{eq:RvMstar} is meaningful because $N^*$
is a named internal counting function. It is not a count obtained by
performing a proper-class contour integral or enumerating surreal zeros
with surreal ordinal multiplicities.

\subsection{Positive proportions and stable infinitesimal zero halos}

Write $N_0(T)$ for the critical-line count with multiplicity, and
$N_0^{\mathrm s}(T)$ for the number of simple critical-line zeros.
A published benchmark of Pratt, Robles, Zaharescu and Zeindler is
\begin{equation}
\liminf_{T\to\infty}\frac{N_0(T)}{N(T)}\geq0.417293,
\qquad
\liminf_{T\to\infty}\frac{N_0^{\mathrm s}(T)}{N(T)}\geq0.407511
\label{eq:PRZZ}
\end{equation}
\cite{PRZZ}. The two constants concern different properties; the first
must not be quoted as a bound for simple critical-line zeros.

The August 2026 preprint of Alp\"oge and Furman, followed by the September
2026 preprint of Lamzouri, gives the stronger reported bound
\begin{equation}
\liminf_{T\to\infty}\frac{N_0^{\mathrm s}(T)}{N(T)}\geq C_0,
\qquad
C_0=\frac32-\frac1{\sqrt2}\cot\frac1{\sqrt2}
       =0.6725007\ldots.
\label{eq:C0}
\end{equation}
Both also give a distinct-zero proportion at least $(1+C_0)/2$;
Lamzouri records additional inequalities for simple-or-critical zeros
\cite{AF,Lamzouri}. These are cited recent preprint results, not arguments
independently verified in this manuscript. Reported formal certificates
in the source projects were not run here. The conclusions below are
parametric in any accepted lower bound $\kappa$, so do not depend on
treating the recent improvement as an independently audited input.

\begin{corollary}[Quantitative robust-halo consequence]
\label{cor:stableproportion}
Suppose a classical theorem gives
$\liminf N_0^{\mathrm s}(T)/N(T)\geq\kappa$.
Then the centre-indexed proportion of uniformly stable critical-line halos
from Section~\ref{sec:robust} has the same lower bound $\kappa$.
In any elementary extension naming the two counting functions, for every
fixed ordinary $c<\kappa$ and every infinite $T$,
\begin{equation}
(N_0^{\mathrm s})^*(T)\geq cN^*(T).
\label{eq:proportionstar}
\end{equation}
\end{corollary}
\begin{proof}
The halo assertion is Theorem~\ref{thm:simplestable} applied to every
simple critical-line centre. For the transferred assertion, the definition
of liminf gives an ordinary $T_c$ such that
$N_0^{\mathrm s}(T)\geq cN(T)$ for every $T\geq T_c$. Transfer preserves
this quantified inequality with the same ordinary $c,T_c$.
\end{proof}

The conclusion is not generally the exact endpoint inequality with
$c=\kappa$. A liminf bound alone does not provide an eventual lower bound
at the endpoint itself. Nor does a positive proportion prove that every
zero is on the line or simple. It does, together with the unbounded total
count, imply the existence of critical-line zeros at infinite heights in
the elementary model.

\subsection{Zero density: transferred bounds and their limitations}

For fixed ordinary $1/2<\sigma<1$, let $N(\sigma,T)$ count zeros with
$\Re\rho\geq\sigma$ and $0<\Im\rho\leq T$. A useful recent benchmark
from Guth and Maynard is
\begin{equation}
N(\sigma,T)\ll T^{A(\sigma)+o(1)},
\qquad A(\sigma)=\frac{15(1-\sigma)}{3+5\sigma}.
\label{eq:GM}
\end{equation}
Combining their estimate with Ingham's gives the convenient uniform-exponent
form
\begin{equation}
N(\sigma,T)\ll T^{\frac{30}{13}(1-\sigma)+o(1)}.
\label{eq:GMsimple}
\end{equation}
Their source counts both signs of the ordinate; these upper bounds therefore
also apply to the upper-half-plane count used here \cite{GM}.
At $\sigma=3/4$, the exponent in \eqref{eq:GM} is $5/9$.

For each fixed ordinary $\sigma$ and ordinary $\eta>0$, an explicit
quantified consequence has the form
\begin{equation}
N^*(\sigma,T)\leq C_{\sigma,\eta}T^{A(\sigma)+\eta}
\label{eq:GMstar}
\end{equation}
for all sufficiently large $T$, including infinite $T$, with an ordinary
constant and with the transported power function. This is transfer of an
ordinary estimate, not a new improvement of its exponent.

Even a bound making $N^*(\sigma,T)/N^*(T)$ infinitesimal does not imply
$N^*(\sigma,T)=0$: an infinite population can have infinitesimal relative
density. Moreover fixed-$\sigma$ bounds do not exclude zeros at
$\Re\rho=1/2+\epsilon$ with positive infinitesimal $\epsilon$. Doing that
would require a theorem whose uniformity actually includes that range,
or RH itself. These are two separate ways in which density information
falls short of a global zero-location statement.

\subsection{Prime counting and finite positivity tests}

The ordinary equivalence between RH and the prime-counting remainder
\begin{equation}
\psi(x)=x+\bigO\!\left(x^{1/2}\log^2x\right)
\label{eq:psiRH}
\end{equation}
uses the classical Chebyshev function and its arithmetic definition
\cite{Titch}. Once $\psi$ and the relevant arithmetic structure are named,
its quantified consequences transfer. In a bare surreal field there is
no specified function called ``the number of primes below $x$'' at an
infinite $x$. Ordinary integers, omnific integers and transferred integers
are different data; writing a sum over unspecified ``surreal primes''
does not define \eqref{eq:psiRH}.

There is also a simple limit on what extending scalar coefficients can
achieve in finite-dimensional positivity criteria.

\begin{proposition}[Finite Hermitian tests do not become stronger]
\label{prop:Hermitian}
Let $A$ be an ordinary complex Hermitian $n\times n$ matrix. It is
positive semidefinite over $\C$ if and only if
\begin{equation}
u^*Au\geq0\qquad\text{for every }u\in\SC^n.
\label{eq:Hermitian}
\end{equation}
\end{proposition}
\begin{proof}
Diagonalize $A=UDU^*$ over $\C$, with ordinary real eigenvalues $d_j$.
The same identity holds over $\SC$. If $d_j\geq0$ for every $j$, then
$u^*Au=\sum_j d_j|(U^*u)_j|^2\geq0$ in the ordered real surreal field.
If some $d_j<0$, its ordinary eigenvector already violates positivity.
\end{proof}

Thus allowing surreal coefficients in a fixed finite family of classical
Weil-type tests adds no condition. A genuinely enlarged infinite test
space requires its own support, summation and positivity theorems. It
cannot be obtained by relabelling the scalars in a finite matrix argument.
Likewise quantities defined by fixed ordinary orders of differentiation
of the canonical lift at ordinary points are exactly the classical
quantities. Adjoining infinitesimals does not alter those coefficients.

\section{Infinitesimal heat time and escape to infinite modulus}
\label{sec:heat}

\subsection{The classical input}

Use the normalization of Rodgers and Tao \cite{RT}:
\begin{align}
H_\tau(z)&=\int_0^\infty e^{\tau u^2}\Phi(u)\cos(zu)\,du,
\label{eq:Hheat}\\
\Phi(u)&=\sum_{n\geq1}
 (2\pi^2n^4e^{9u}-3\pi n^2e^{5u})e^{-\pi n^2e^{4u}},
\notag\\
H_0(z)&=\frac18\xi\!\left(\frac12+\frac{iz}{2}\right).
\label{eq:H0}
\end{align}
The de Bruijn--Newman constant $\Lambda$ is characterized by: all zeros
of $H_\tau$ are real exactly when $\tau\geq\Lambda$. Rodgers and Tao
prove $\Lambda\geq0$. Consequently every ordinary negative $\tau$ has
at least one nonreal zero of $H_\tau$.

The super-exponential decay of $\Phi$ makes $H_\tau(z)$ jointly analytic
in complex $(\tau,z)$ on bounded parameter sets, by differentiation under
the ordinary integral. Indeed the factors from differentiating are
polynomials in $u$, multiplied by $e^{Cu^2+Cu}$ on a fixed compact set,
and remain dominated by the super-exponential tail. For real $\tau$ the
function is conjugation-symmetric, and direct differentiation gives
$\partial_\tau H_\tau=-\partial_z^2H_\tau$.

No surreal integral is being used here: these are ordinary analytic germs
before lifting or transfer.

\begin{lemma}[Local real-root persistence]
\label{lem:heatpersist}
If $a$ is a simple real zero of $H_0$, then there are ordinary $r,\delta>0$
such that for every real $|\tau|<\delta$, all zeros of $H_\tau$ in
$|z-a|<r$ are real and there is exactly one such zero. Its root germ has
\begin{equation}
z(\tau)=a+\tau\frac{H_0''(a)}{H_0'(a)}+\bigO(\tau^2).
\label{eq:heatroot}
\end{equation}
The same local uniqueness and reality statements hold after canonical
infinitesimal evaluation or elementary transfer.
\end{lemma}
\begin{proof}
The analytic implicit-function theorem at $(0,a)$ gives one local root
and a unit factorization. Shrink ordinary $r,\delta$ so that it is the
only root in the disk. For real $\tau$, conjugation preserves this disk
and the function; uniqueness makes the root real. Differentiating
$H_\tau(z(\tau))=0$ gives
$z'(0)=-\partial_\tau H/H_z=H_0''(a)/H_0'(a)$.
Hahn evaluation preserves the factorization. Alternatively the existence,
uniqueness and reality statement with its ordinary radii is first order
in the real-coordinate functions $H$, so transfers.
\end{proof}

\begin{theorem}[Conditional escape at negative infinitesimal time]
\label{thm:escape}
Include the real-coordinate functions $H_\tau(z)$ in an elementary
extension as in Section~\ref{sec:transfer}. Assume the ordinary RH and
simplicity of every nontrivial zeta zero. For every negative infinitesimal
$\tau\in F$, the transferred function $H^*_{\tau}$ has a nonreal zero,
but every such nonreal zero has infinite modulus.

By contrast, the canonical joint Taylor lift at $(0,\C)$, evaluated at
that infinitesimal parameter, has only real zeros on its finite $z$-plane.
\end{theorem}
\begin{proof}
The unconditional ordinary statement
\[
\forall\tau<0\ \exists z\quad
\bigl(H_\tau(z)=0\text{ and }\Im z\ne0\bigr)
\]
follows from $\Lambda\geq0$ and the defining characterization of
$\Lambda$. Transfer supplies a nonreal zero for every negative
infinitesimal $\tau$.

Suppose such a zero $z$ had finite modulus, and let $a=\st z$.
Ordinary continuity of the joint function at $(0,a)$, applied with each
ordinary positive error tolerance and then transferred, gives
$\st(H^*_{\tau}(z))=H_0(a)$. Thus $H_0(a)=0$.
By \eqref{eq:H0}, RH and simplicity make $a$ a simple real zero.
Lemma~\ref{lem:heatpersist}, transferred with its ordinary radii, then
forces $z$ real, a contradiction. Therefore every nonreal zero has
infinite modulus.

For the canonical joint lift the same standard-part argument puts every
finite zero in the halo of a zero of $H_0$. The evaluated unit
factorization in Lemma~\ref{lem:heatpersist} makes the unique root in
that halo real. There is no domain of infinite $z$ in this finite-only
construction, so it has no nonreal zeros at all on its stated domain.
\end{proof}

This is a concrete explanation of why infinitesimal stability at every
ordinary zero does not contradict the criticality expressed by
$\Lambda\geq0$. Under RH and simplicity, the instability is invisible
on every fixed ordinary bounded region and is witnessed only at infinite
modulus in the transferred model. The assertion concerns modulus; no
specific asymptotic scale for the first nonreal zero is proved here.

Without the two conjectural assumptions, the proof still gives an
unconditional localization statement: a finite nonreal zero at infinitesimal
time must have standard part either a nonreal zero of $H_0$ or a multiple
real zero of $H_0$. It cannot have a simple real zero as its standard part.
This is a useful diagnostic that does not presuppose RH.

\section{Synthesis and concrete further questions}
\label{sec:synthesis}

The canonical finite theory is fully specified on its domain:
Gamma and zeta are defined, their ordinary identities lift, their divisors
are exactly known in terms of the classical divisors, and its RH is
precisely classical RH. The robust version is strictly more informative
as a formulation: it packages RH and simplicity into stability under
symmetric infinitesimal perturbations.

The infinite-right theory supplies a second exact construction. Strong
Dirichlet summability, rather than ordinary convergence, gives an injective
realization of all ordinary arithmetic coefficient functions. Euler products,
M\"obius inversion, logarithmic derivatives and zero-freeness are literal
identities. The infinite-shift Hurwitz construction provides a compatible
formal bridge to the Stirling logarithm, while explicitly avoiding the
nonsummable shifted ordinary series.

At infinite height, the global problem requires additional choices.
The canonical published phase $E_0$ is a legitimate exponential but
forgets the frequencies needed by classical Dirichlet oscillation. The
resonance theorem proves that certain natural continuation and divisor
requirements cannot all coexist with it. Weak local analytic and symmetry
axioms leave the global zero set undetermined, while elementary extension
gives a coherent alternative on embedded set-sized models, with its own
transported operations.

Several sharply formulated projects emerge from these results.

\paragraph{Nonresonant global completion.}
Specify a phase $U$ satisfying at least the two-frequency test
\eqref{eq:twofreq}, a summation rule in a finite-real-part right half-plane
at infinite height, and a meromorphic continuation class. Determine whether
one can simultaneously retain a single zeta pole, the completed functional
equation and a prescribed Gamma asymptotic normalization. The theorem here
identifies a necessary obstruction to avoid, not a sufficient construction.

\paragraph{Compatibility of canonical and transferred local values.}
An elementary model has exactly the correct finite zeros, but an arbitrary
ordered-field embedding does not identify all its values with Hahn lifts.
Find additional hypotheses on an embedding or on a restricted analytic
substructure that enforce such equality, and determine how far this
compatibility extends into infinite sectors. Restricted-analytic surreal
extension results in \cite{CE,EK} provide relevant structure, but do not
by themselves settle a global zeta-and-phase compatibility statement.

\paragraph{Quantitative heat escape.}
Under RH and simplicity, define in a transferred model the least scale, or
an appropriate infimum scale, at which nonreal zeros appear for
$\tau=-\epsilon$. Theorem~\ref{thm:escape} proves only that every ordinary
bound is exceeded. Relating that scale to classical small gaps and the
heat-flow collision mechanism would turn the qualitative infinite-modulus
statement into asymptotic information.

\paragraph{Certified local algebra.}
The divisor theorem, simple-root symmetry argument, splitting-scale
calculation and finite Dirichlet convolution identities have clean
formalization boundaries. They can be formalized without assuming RH.
The support and class-size infrastructure, and any imported analytic
number-theory theorem, must be recorded separately rather than hidden
inside a computational certificate.

\medskip
There is therefore no single unqualified statement called ``RH over the
surcomplex numbers.'' There are precise statements attached to precise
extensions. Two of them---the canonical finite RH and the elementary
RH---are equivalent to ordinary RH for different reasons. A third, robust
finite RH, is equivalent to RH plus simplicity. For a separately chosen
global Hahn-analytic or phase-dependent completion, RH is a genuinely
additional assertion whose content is inseparable from the construction.

\appendix
\section{Proof dependencies and reproducible finite checks}
\label{app:checks}

The supplied Python program passed 221 exact symbolic checks
(Python 3.13.5, SymPy 1.14.0). It performs checks in ordinary
polynomial and truncated formal-series rings. It checks the Stirling shift
coefficients, the Hurwitz shift expansion for several ordinary parameters,
Bernoulli special values, the derivative bridge at zero, the first two
implicit-root coefficients, finite Dirichlet-convolution identities, and
finite Euler-product coefficients. It also verifies the algebra of the
symmetric quartet construction and several local splitting polynomials.

These checks detect finite algebraic errors. They do not prove Hahn
summability, quantify over surreal classes, verify analytic continuation,
prove the classical RH, or validate any of the cited analytic-number-theory
papers. No arbitrary large floating-point number is substituted for an
infinite surreal. The general arguments are the proofs in the text.

\begin{center}
\begin{tabular}{P{0.34\textwidth}P{0.54\textwidth}}
\toprule
Result & Principal dependencies\\
\midrule
Finite divisor preservation & Neumann support lemma; ordinary meromorphic local factorization\\[3pt]
Robust RH equivalence & Finite divisor preservation; ordinary analytic implicit-function theorem; xi symmetries\\[3pt]
Infinite-right Dirichlet algebra & Purely-infinite exponential monomials; Neumann lemma; unique prime factorization\\[3pt]
Hurwitz--Stirling bridge & Bernoulli generating identity; formal shift calculus; infinitesimal evaluation\\[3pt]
Phase obstruction & Published phase-forgetting exponential; ordinary meromorphic identity theorem; classical zeta values\\[3pt]
Elementary RH and estimates & Elementary extension; ordered-field embedding; explicitly named classical functions and estimates\\[3pt]
Heat escape & Joint analytic root persistence; Rodgers--Tao nonnegativity; RH and simplicity for the conditional conclusion\\
\bottomrule
\end{tabular}
\end{center}

The source is self-contained LaTeX with an internal bibliography. Running
\texttt{./build.sh} compiles the article; running
\texttt{python code/checks.py} reruns the finite checks and prints a JSON
report. An optional explicit output path writes a new report; the default
does not overwrite the shipped evidence. See the accompanying README for
file descriptions and the recorded build status.

\section{Notation and domain index}

\begin{longtable}{P{0.22\textwidth}P{0.66\textwidth}}
\toprule
Symbol & Meaning and scope\\
\midrule
\endhead
$\SC=\No[i]$ & Surcomplex field; class-sized globally, with set-sized supports in every calculation.\\[3pt]
$\OO$, $\mm$ & Finite surcomplex elements and infinitesimals.\\[3pt]
$\st z$ & Ordinary complex standard part of a finite element.\\[3pt]
$\pp x$, $\fin x$ & Purely infinite and finite normal-form parts of a real surreal.\\[3pt]
$c+\mm$ & Infinitesimal halo of an ordinary centre.\\[3pt]
$b+\OO$ & Finite-width galaxy about a possibly infinite point.\\[3pt]
$\Sh$, $\Ph$ & Strong Hahn sum and product defined by strongly summable expansion.\\[3pt]
$f^\#$ & Canonical finite Taylor--Laurent lift of an ordinary function.\\[3pt]
$\xi$ & Completed zeta in normalization \eqref{eq:xi}; all its zeros are nontrivial zeta zeros.\\[3pt]
$F_{\epsilon,h}$ & Symmetric infinitesimal perturbation of the finite xi lift.\\[3pt]
$\mathcal D_+$, $\zD$ & Positive infinite real part, finite imaginary part; exact strong-Dirichlet zeta.\\[3pt]
$\LSt$, $\Gamma_{E,\mathrm{St}}$ & Strong Stirling logarithm and its exponentiation with a specified global phase.\\[3pt]
$Z_\infty(s,A)$ & Infinite-positive-shift Hurwitz prescription; ordinary real parameter and restricted parameter germs.\\[3pt]
$E_0$ & Published phase-forgetting surcomplex exponential with kernel $2\pi i\Oz$.\\[3pt]
$f^*$ & Named function in an elementary extension; not automatically a canonical Hahn lift after embedding.\\[3pt]
$N^*(T)$ & Transferred internal zero count, not an external cardinality.\\[3pt]
$H_\tau$ & Ordinary de Bruijn--Newman heat family before lifting or transfer.\\
\bottomrule
\end{longtable}

\begin{thebibliography}{99}
\small
\raggedright
\interlinepenalty=10000
\setlength{\itemsep}{5pt}
\setlength{\parskip}{0pt}
\addcontentsline{toc}{section}{References}

\bibitem{Conway}
J. H. Conway, \emph{On Numbers and Games}, second edition, A K Peters, 2001.

\bibitem{Gonshor}
H. Gonshor, \emph{An Introduction to the Theory of Surreal Numbers},
London Mathematical Society Lecture Note Series 110, Cambridge University
Press, 1986.

\bibitem{CE}
O. Costin and P. Ehrlich, \emph{Integration on the surreals},
arXiv:2208.14331v5. In particular the finite analytic extension framework,
the Gamma example, and Proposition 93.
\url{https://arxiv.org/html/2208.14331v5}.

\bibitem{EK}
P. Ehrlich and E. Kaplan, \emph{Surreal ordered exponential fields},
arXiv:2002.07739v3. Section 11, especially Proposition 11.2 and Theorem 11.1,
for surcomplex exponential extensions and integer parts.
\url{https://arxiv.org/abs/2002.07739v3}.

\bibitem{repo}
V. Reshetnikov, \emph{Surreal} repository, documentation catalogue,
pinned snapshot \pin. AI-assisted research collection; targeted inspection
of the maintained documentation guides.
\href{https://github.com/VladimirReshetnikov/Surreal/tree/905dd19954db43ac81f76f72036520d0eadc5710/docs}{Pinned documentation catalogue}.

\bibitem{repoGamma}
\repo, \emph{Gamma functions over the surreals: a sharp convexity
classification}, maintained report guide and linked article, directory
\texttt{docs/surreal/gamma-functions}, at the snapshot in \cite{repo}.
\href{https://github.com/VladimirReshetnikov/Surreal/tree/905dd19954db43ac81f76f72036520d0eadc5710/docs/surreal/gamma-functions}{Pinned Gamma report directory}.

\bibitem{repoAnalysis}
\repo, \emph{Surcomplex analysis: holomorphic functions on
$K=\No[i]$}, maintained report guide, directory
\texttt{docs/surcomplex/analysis}, at the snapshot in \cite{repo}.
\href{https://github.com/VladimirReshetnikov/Surreal/tree/905dd19954db43ac81f76f72036520d0eadc5710/docs/surcomplex/analysis}{Pinned analysis report directory}.

\bibitem{repoTrig}
\repo, \emph{Trigonometry on the Surcomplex Plane}, maintained report
guide, directory \texttt{docs/surcomplex/trigonometry}, at the snapshot
in \cite{repo}.
\href{https://github.com/VladimirReshetnikov/Surreal/tree/905dd19954db43ac81f76f72036520d0eadc5710/docs/surcomplex/trigonometry}{Pinned trigonometry report directory}.

\bibitem{DLMFgamma}
NIST Digital Library of Mathematical Functions, Section 5.2,
\emph{Definitions: Gamma and psi functions}.
\url{https://dlmf.nist.gov/5.2}.

\bibitem{DLMFidentities}
NIST Digital Library of Mathematical Functions, Section 5.5,
\emph{Functional relations for the Gamma function}.
\url{https://dlmf.nist.gov/5.5}.

\bibitem{DLMFstirling}
NIST Digital Library of Mathematical Functions, Section 5.11,
\emph{Asymptotic expansions for the Gamma function}.
\url{https://dlmf.nist.gov/5.11}.

\bibitem{DLMFzeta}
NIST Digital Library of Mathematical Functions, Section 25.2,
\emph{Definition and expansions of the Riemann zeta function}.
\url{https://dlmf.nist.gov/25.2}.

\bibitem{DLMFfe}
NIST Digital Library of Mathematical Functions, Section 25.4,
\emph{Reflection formulas for the Riemann zeta function}.
\url{https://dlmf.nist.gov/25.4}.

\bibitem{DLMFzeros}
NIST Digital Library of Mathematical Functions, Section 25.10,
\emph{Zeros of the Riemann zeta function}.
\url{https://dlmf.nist.gov/25.10}.

\bibitem{DLMFhurwitz}
NIST Digital Library of Mathematical Functions, Section 25.11,
\emph{Hurwitz zeta function}. Especially recurrence, Bernoulli special
values, parameter derivative at zero and large-parameter asymptotics.
\url{https://dlmf.nist.gov/25.11}.

\bibitem{Titch}
E. C. Titchmarsh, \emph{The Theory of the Riemann Zeta-Function}, second
edition revised by D. R. Heath-Brown, Oxford University Press, 1986.

\bibitem{PT}
D. Platt and T. Trudgian, \emph{The Riemann hypothesis is true up to
$3\cdot10^{12}$}, Bulletin of the London Mathematical Society (2021),
DOI: 10.1112/blms.12460; arXiv:2004.09765.
\url{https://arxiv.org/abs/2004.09765}.

\bibitem{MTY}
M. J. Mossinghoff, T. S. Trudgian and A. Yang,
\emph{Explicit zero-free regions for the Riemann zeta-function},
arXiv:2212.06867, 2022.
\url{https://arxiv.org/abs/2212.06867}.

\bibitem{PRZZ}
K. Pratt, N. Robles, A. Zaharescu and D. Zeindler,
\emph{More than five-twelfths of the zeros of $\zeta$ are on the critical
line}, Research in the Mathematical Sciences \textbf{7}, article 2 (2020);
arXiv:1802.10521v3. DOI: 10.1007/s40687-019-0199-8.
\url{https://arxiv.org/abs/1802.10521v3}.

\bibitem{GM}
L. Guth and J. Maynard, \emph{New large value estimates for Dirichlet
polynomials}, arXiv:2405.20552, 2024. Theorem 1.2 and equation (1.4)
for the zero-density bounds used here.
\url{https://arxiv.org/abs/2405.20552}.

\bibitem{AF}
L. Alp\"oge and R. Furman, \emph{More than two thirds of the zeta zeros
are simple and on the critical line}, arXiv:2608.13637v2,
19 August 2026. Recent preprint; cited as an external result, not
independently verified here.
\url{https://arxiv.org/abs/2608.13637v2}.

\bibitem{Lamzouri}
Y. Lamzouri, \emph{A new proof that more than $2/3$ of the zeros of the
Riemann zeta function are simple and on the critical line},
arXiv:2609.02882v2, 8 September 2026. Theorem 1.1.
Recent preprint; cited as an external result, not independently verified here.
\url{https://arxiv.org/abs/2609.02882v2}.

\bibitem{RT}
B. Rodgers and T. Tao, \emph{The de Bruijn--Newman constant is
non-negative}, arXiv:1801.05914v5, corrected version of the published
paper, 3 July 2021. The normalization of $H_0$ is equation (1), and
nonnegativity is Theorem 1.
\url{https://arxiv.org/html/1801.05914v5}.

\end{thebibliography}
\end{document}
